\section{Introduction}\label{sec:introduction}

Many polynomial-time algorithms make intermediate choices without changing
their final answer. Gaussian elimination, for example, may select different
pivots while deciding the same consistency question. The difficulty on an
unordered finite structure is to permit intermediate choices in a logic
whose answers are invariant under renaming the input elements. Witnessed symmetric
choice addresses this difficulty by requiring certificates for the
symmetry of every choice set~\cite[Introduction]{ls2023}.

A computation on a finite relational structure is \emph{choiceless} when it
does not select an element from an unordered set merely because that set is
nonempty. Choiceless polynomial time with counting, denoted by \(\CPT\),
performs invariant computations with hereditarily finite sets over the input
atoms. It can form all objects in a definable finite family and count their
members, but it has no arbitrary ordering of the atoms. Witnessed symmetric
choice allows a more specific operation: select an element of an orbit while
constructing automorphisms that certify the permitted symmetry.

We use exactly the counting version of \(\CPT\) and the extension
\(\CPT+\WSC\) in Sections~2 and~3 of Lichter and Schweitzer's
version~3~\cite{ls2023}. The witnessing automorphisms preserve
the entire input vocabulary and stabilize both the parameters and the
preceding intermediate states. Section~\ref{sec:formalism} specifies the
evaluation rules and polynomial resource bounds. We use \(\CPT\) for the
counting model throughout; some of the older literature reserves that
notation for the model without counting.

For a finite vocabulary \(\tau\), a Boolean query is an isomorphism-invariant
class of finite \(\tau\)-structures. A \(\CPT+\WSC\) sentence may return
true, false, or a failure value \(\dagger\). Its \emph{true-model class}
consists of the structures on which its value is true; false and
\(\dagger\) both lie outside that class. Our separating sentence always
returns a Boolean value, and satisfies the following strict
expressiveness result.

\Needspace{10\baselineskip}
\begin{theorem}\label{thm:main}
There exist a vocabulary \(\tau\) of eight binary relations, a fixed polynomial
\(p\), and a \(\CPT+\WSC\) sentence \((\Phi,p)\) with one occurrence of a WSC
operator such that:
\begin{enumerate}[label=\textup{(\roman*)}]
\item \((\Phi,p)\) has a Boolean value on every finite \(\tau\)-structure;
\item no \(\CPT\) sentence over \(\tau\) has the same true-model class.
\end{enumerate}
The \(\CPT\) sentences in \textup{(ii)} may use counting and hereditarily
finite sets of unrestricted finite rank.
\end{theorem}

Since every \(\CPT\) sentence is already a \(\CPT+\WSC\) sentence, the theorem
gives a strict inclusion of their collections of definable Boolean queries.
In the true-model sense above, this answers affirmatively the expressiveness
question posed in Section~7 of~\cite{ls2023}.
One syntactic occurrence can make many successive choices during its
iteration and can be evaluated at different parameter tuples.
The result concerns strict expressive power; capture of all polynomial-time
queries and the expressive effect of nested WSC operators remain separate
questions.

\paragraph{Background and relation to earlier work.}
Blass, Gurevich and Shelah introduced choiceless polynomial-time computation
using hereditarily finite sets and proposed its extension by a cardinality
operation~\cite{bgs1999}. Their subsequent study of unordered computation
examined Cai--F{\"u}rer--Immerman structures and finite-field linear algebra
as possible sources of separation~\cite[Sections~4 and~6]{bgs2002}.
Shelah also made an early noncapture claim for counting extensions of
choiceless computation~\cite[Introduction and Section~4]{shelah2000}.
Later accounts continued to formulate the counting capture problem as
open~\cite[Section~7]{bgs2002}\cite[Introduction]{pago2023}.
We retain this historical distinction and make no priority claim for
noncapture of the counting formalism.

The construction of Cai, F{\"u}rer and Immerman hides a global parity
obstruction from bounded-variable counting logic~\cite{cfi1992}.
Its local constraints and movable obstruction are predecessors of the
charged structures used here. Passing from counting logic to full CPT
requires more. Dawar, Richerby and Rossman showed that the preordered CFI
query arising from ordered base graphs is definable without counting,
but cannot be computed even with counting when the rank of activated
sets is restricted to \(o(\log n/\log\log n)\)
\cite[Theorems~17 and~40, author version]{drr2008}.
Pakusa, Schalth{\"o}fer and Selman extended positive CPT definability
results to CFI problems over preordered base graphs with logarithmic
color classes and, separately, over base graphs whose maximum degree is
at least a fixed positive fraction of their order
\cite[Theorems~1 and~15]{pakusaSchalthoeferSelman2016}.
These results show why a lower bound for the full formalism must account
for its ability to construct nested sets.

Positive CPT results also exploit local automorphism groups. Suppose a
total preorder is supplied on the input, with every equivalence class of
size at most a fixed constant. Abu Zaid, Gr{\"a}del, Grohe and Pakusa gave
CPT canonization when the substructure induced by each class has an
abelian automorphism group~\cite[Definition~15 and Corollary~19]{abuZaid2014}.
Canonization assigns an isomorphic structure on an ordered universe in an
isomorphism-invariant way. Lichter and Schweitzer extended such
canonization results to structures of arity at most three whose class groups
are dihedral groups of odd degree or cyclic groups, and to graphs with
dihedral or cyclic class groups
\cite[Theorem~1]{lichterSchweitzer2021}.
These hypotheses constrain the automorphism groups of the induced class
substructures in addition to bounding the class sizes.

Other lower bounds address different parts of this difficulty. Rossman
ruled out CPT construction of the set of hyperplanes of an input vector
space over a fixed finite field
\cite[Theorem~6.1, author version]{rossman2010}.
Pago ruled out sufficiently fine total preorders of hypercubes
\cite[Theorem~2 and Corollary~3, full version]{pago2021} and developed
a symmetric-XOR-circuit approach to restricted classes of CPT algorithms
\cite[Theorems~1--2, full version]{pago2023}.
The first two are functional obstructions; the circuit translation and
lower bound have different restrictions. None of these statements alone
is a lower bound for an arbitrary full-CPT Boolean query.

Symmetry-restricted choice has a separate history. Gire and Hoang
introduced an efficiently evaluated construct based on specified
symmetries~\cite{gireHoang1998}. Dawar and Richerby developed a fixed-point
logic with symmetric choice, including parameters and nesting
\cite{dawarRicherby2003}. Lichter and Schweitzer's extension places
witnessed choice within the HF model and permits witnesses to inspect
the terminal state while requiring them to preserve the preceding
history~\cite{ls2023}. Their Theorem~1 turns definable isomorphism testing
into canonization on classes closed under individualization. This
conditional result explains the power of the extension without assuming
that it captures polynomial time on all structures.
Lichter's subsequent separations and results on nesting witnessed-choice
and interpretation operators concern extensions of fixed-point logic
with counting; they use that smaller base logic
\cite[Theorems~1 and~3, version~8]{lichter2026}.

The charged structures and their support and counting comparison were
developed in the companion \emph{Choiceless polynomial time with counting
does not capture polynomial time}\linebreak
\cite[Sections~2--5]{cptCompanion2026}.
We give the complete arguments needed here. In particular, we retain its
rank-independent support estimate, base equivalence, quantitative
homogeneity and supported-HF transfer
\cite[Propositions~8, 11 and~13, Theorem~14]{cptCompanion2026}.
The companion's noncapture theorem concerns a linear-consistency query;
it does not itself supply a definition in the witnessed-choice language.
Our all-input query is specified separately below. The additional steps
are its legal witnessed-choice definition and a direct comparison of
arbitrary CPT terms on the common grid pairs. The latter does not use
the interpretation-program characterization employed in the companion.

The support-and-transfer method also has its own ancestry. Blass,
Gurevich and Shelah introduced representations by forms evaluated at
supporting tuples, called molecules
\cite[Sections~8--9, published version]{bgs1999}.
Dawar, Richerby and Rossman developed the counting transfer under
hereditary-support hypotheses
\cite[Section~8, Theorem~33, author version]{drr2008}.
That transfer already permits arbitrary finite rank. The rank limitation
in their lower bound comes from obtaining the required supports. Here
the nonabelian group action forces short supports for its central
subgroup independently of rank, and the quantitative transfer compares
all objects with such hereditary supports.

\paragraph{Proof strategy.}
The input vocabulary describes two sorts of atoms. An \emph{edge atom}
specifies a state on an edge of a three-dimensional box, and a
\emph{configuration atom} specifies compatible states on the edges incident
with a vertex. State vectors agree across the sides of a square face;
scalar states satisfy a prescribed divergence at each vertex. For each
box size we construct two structures of equal cardinality: one has total
divergence zero and the other has total divergence one.
Globally consistent edge and vertex choices exist in
the zero-charge structure and not in the unit-charge structure.

The witnessed-choice sentence first uses a tuple of four vertex blocks to
define an order of the underlying graph and a spanning tree. It then
chooses states only on edges outside that tree. Two candidates at the next
edge are connected by two explicit transformations. A translation on at
most four faces aligns their vector coordinates and moves only a fixed
number of atoms. A circulation around the new edge and its tree return
path corrects the remaining scalar difference. Neither transformation
disturbs previously selected edge states. The formula enumerates the
possible local maps and tree-cycle maps, and keeps only composites that
are automorphisms of the full input fixing the required history. A
transitivity guard and a pure dummy choice make the same sentence legal
on arbitrary structures.
The cycle-witness method of Gire and Hoang~\cite{gireHoang1998} is
reconstructed in \cite[Section~6]{ls2023}, including preservation of
earlier choices.
Here the cycle circulation follows a bounded local vector correction;
the resulting maps are checked against the full input and its history
(Proposition~\ref{grid:transitivity}).

After the choices, a monotone elimination on the tree decides whether the
selected edge states extend to compatible configurations. The zero-charge
inputs always pass; the unit-charge inputs always fail.
Sections~\ref{sec:structures}--\ref{sec:grid-behavior} develop this
construction and its resource bound.

The lower bound begins with a central subgroup \(K\) of the automorphism
group. An HF object has a \emph{\(K\)-support} if fixing a specified tuple
of atoms pointwise forces every element of \(K\) to fix that object.
A small orbit under a larger group forces a short \(K\)-support.
Counting equivalence and a quantitative homogeneity property of the base
structures then transfer to the entire domains of hereditarily supported
HF objects. These domains have no rank bound.
Sections~\ref{sec:support}--\ref{sec:hf-transfer} prove those statements.

Finally, for each fixed \(\CPT\) sentence, the complete actual evaluation
trace has polynomial size and is invariant. Its values therefore lie in
the supported domains. We describe its evaluation by a counting formula
using a fixed number of variable names. The description may grow in
length with the common input size, but its variable width does not grow.
The crucial induction makes every term graph unique over all proposed
outputs whenever its arguments are an actual activation. This prevents
arbitrary high-rank objects from producing spurious fixed points or
altering an empty-set default. Section~\ref{sec:cpt-lower-bound} applies
this description to force equal \(\CPT\) answers on sufficiently large
pairs, completing the separation.
